\documentclass{article}

\usepackage{booktabs}
\usepackage{tabularx}
\usepackage{makecell}

\title{Development Plan\\\progname}

\author{\authname}

\date{}

\input{../Comments}
%% Common Parts

\newcommand{\progname}{Project Proxi} % PUT YOUR PROGRAM NAME HERE
\newcommand{\authname}{Team \#23, Project Proxi
\\ Savinay Chhabra
\\ Amanbeer Singh Minhas
\\ Gourob Podder
\\ Ajay Singh Grewal} % AUTHOR NAMES                  

\usepackage{hyperref}
    \hypersetup{colorlinks=true, linkcolor=blue, citecolor=blue, filecolor=blue,
                urlcolor=blue, unicode=false}
    \urlstyle{same}
                                


\begin{document}

\maketitle

\begin{table}[hp]
\caption{Revision History} \label{TblRevisionHistory}
\begin{tabularx}{\textwidth}{llX}
\toprule
\textbf{Date} & \textbf{Developer(s)} & \textbf{Change}\\
\midrule
September 21, 2025 & Savinay Chhabra  & Created Draft Development Plan; Sections 1, 2, 3, 4, 5, 7\\
September 22, 2025 & Savinay Chhabra  & Sections 6, 8\\ 
September 22, 2025 & Ajay Singh Grewal & Section 9, 10, 11\\
September 22, 2025 & Amanbeer Singh Minhas & Team Charter and Reflection\\
April 6, 2026 & Ajay Singh Grewal & Final docs, implemented TA feedback\\
\bottomrule
\end{tabularx}
\end{table}

\newpage{}
This report outlines the development plan for the Proxi AI Assistant. 
The document details the following sections of the Development Plan:
\begin{itemize}
  \item Confidentiality Information
  \item IP Protection
  \item Copyright License
  \item Team Meeting Plans
  \item Team Communication Plans
  \item Team Member Roles
  \item Workflow Plan
  \item Project Decomposition and Scheduling
  \item Proof of Concept Demo Plan
  \item Expected Technology
  \item Coding Standard
\end{itemize}

\section{Confidential Information?}

There is no confidential information to protect.

\section{IP to Protect}

There is no IP to protect.

\section{Copyright License}

This project adopts the MIT License. LICENSE file can be found \href{https://github.com/gkpodder/Capstone/blob/main/LICENSE}
{here}. The main project repository can be found \href{https://github.com/gkpodder/Proxi}{here}.

\section{Team Meeting Plan}
The Team will meet every Wednesday from 4:00PM - 4:45PM in person. Campus library meeting rooms will be booked weekly, depending on availability.
\begin{itemize}
  \item A rotating chair will lead each meeting and prepare the meeting\'s agenda in advance.
  \item A rotating notetaker will record meeting minutes, action items and decisions made in the meeting.
  \item Feature requests and Decisions will be documented using GitHub Issues
  \item Meetings with graduate students speacializing in Human-Computer Interfaces will occur biweekly. These will either be virtual or in-person depending on availability.
\end{itemize}

\section{Team Communication Plan}

Urgent issues or clarifications will be communicated through Discord. Discord will also be used for any virtual meetings between the team. 
Meetings with industry advisors and graduate students will take place over MS Teams if held virtually. 
GitHub issues will the the source of truth for techincal information and progress tracking.

\section{Team Member Roles}

\begin{table}[!h]
\caption{Team Member Roles} \label{TblMemberRoles}
\begin{tabularx}{\textwidth}{lXl}
\toprule
\textbf{Role} & \textbf{Responsibilities} & \textbf{Member(s}\\
\midrule
Meeting Chair & Sets meeting agenda and leads weekly meetings & Rotates weekly\\
Note Taker & Records meeting minutes and decisions & Rotates weekly\\
Team Liaison & Primary point of contact for project supervision and industry advisors & Savinay Chhabra \\
Subject Matter Expert & Develops deep knowledge of the project\'s domain & Gourob Podder \\
Test Lead & Responsible for planning, organizing and overseeing testing process. & Amanbeer Singh Minhas \\
Reviewer & Responsible for evaluating Pull Requests to ensure that they adher to team standards. & Ajay Singh Grewal \\
\bottomrule
\end{tabularx}
\end{table}

Role assignments are intended to reflect each member's strengths and areas of domain or technical expertise at the start of the project; however, they are not fixed. If a team member's availability changes, or if the project's needs shift (e.g., increased testing load, a new integration requiring specialist knowledge), roles will be reassigned through team consensus at a regular weekly meeting. Any role change will be recorded in the meeting minutes and reflected in an updated version of this table. The rotating Meeting Chair and Note Taker roles already demonstrate this flexibility in practice, and the same principle applies to all other roles.

\section{Workflow Plan}

\begin{itemize}
  \item Each Feature/Bug/Defect will have a GitHub Issue created for it. This issue will included detailed information along with current progress, assignee and appropriate labels. GitHub Projects will be used to track all incoming and in-progress issues.
	\item Each Issue will a branch for any associated development work to be done.
	\item Once the author is satisfied with their solution, they can put up a Pull Request to merge their changes into the main branch. Each Pull request must link to one or more issues and have a detailed description of the changes made and the user impact of those changes.
	\item After the Pull Request has been raised and all required pre-submit checks have succeded, the author can add the ``Review Needed'' label to their PR.
	\item Each PR must have at least 2 reviewers. After the reviewers have approved, the PR can be merged and the associated issues can be closed.
\end{itemize}

\section{Project Decomposition and Scheduling}

The team will use the Kanban Board on GitHub Projects to track incoming and in-progress issues.
The Board for the project can be found \href{https://github.com/users/gkpodder/projects/2}{here}.


\begin{table}[!h]
\caption{Project Scheduling} \label{ProjScheduling}
\begin{tabularx}{\textwidth}{lXl}
\toprule
\textbf{Deliverable} & \textbf{Deadline} \\
\midrule
Problem Statement, POC Plan, Development Plan & September 22, 2025\\
Req. Doc. and Hazard Analysis Revision 0 & October 6, 2025 \\
V\&V Plan Revision 0 & October 27, 2025 \\
Design Document Revision -1 & November 10, 2025 \\
Design Document Revision 0 & January 19, 2026 \\
Revision 0 Demonstration &  February 13, 2026\\
V\&V Report and Extras Revision 0 & March 09, 2026 \\
Final Demonstration (Revision 1)  &  March 29, 2026 \\
EXPO Demonstration & TBD \\
Final Documentation (Revision 1) & April 06, 2026 \\
\bottomrule
\end{tabularx}
\end{table}

\section{Proof of Concept Demonstration Plan}

What is the main risk, or risks, for the success of your project?  What will you
demonstrate during your proof of concept demonstration to convince yourself that
you will be able to overcome this risk?

The main feasibility risk for Proxi is whether an AI assistant can reliably 
control a user's computer well enough to complete a meaningful task from a 
natural-language command. Since the core value of Proxi is accessible, 
hands-free computer interaction, the project depends on proving that the 
assistant can correctly interpret a request, understand the user's environment, 
and execute the proper action on the machine. Concerns such as security, 
hallucinations, and long-term user trust are also important, but for this proof 
of concept the primary question is whether the basic interaction is technically 
feasible at all.

Our November proof of concept demonstration will therefore focus on one concrete 
end-to-end workflow. A user will issue a voice or text command such as opening 
an application, locating a file, summarizing its contents, or creating a 
calendar event, and Proxi will interpret the intent, use the required tools, 
and visibly complete the task. This will demonstrate that the team's chosen 
automation technology, language model pipeline, and system architecture are 
sufficient to support real computer interaction rather than only simulated 
responses.

If this proof of concept fails, it would indicate that the current approach to 
AI-driven computer control is too unreliable or too difficult to implement 
within the Capstone timeline. In that case, the project would need to reduce 
its scope toward a narrower, more supervised assistant that supports a smaller 
set of predefined tasks instead of broader autonomous computer control.


\section{Expected Technology}

\begin{itemize}
  \item Programming Language: Python for the AI agent, tool plugins, and overall backend structure. TypeScript and React for desktop UI
  \item Agent Core Libraries: pydantic, asyncio, httpx
  \item System Automation: For Windows we will use pywin32, UIAutomationCore; for macOS we will use pyobjc, some AppleScript/JXA bridge; For Browser we will use Playwright
  \item Storage/Database: SQL
  \item Pre-trained Models: OpenAI GPT-4o-mini API for natural language understanding and generation
  \item Local LLM For Future (if needed)
  \item Linter Tools: pylint, Prettier, ESLint
  \item CI/CD: Git, GitHub
\end{itemize}
\section{Coding Standard}

The team will strictly adhere to and follow the following coding standards to ensure clarity, maintainablity and readability:
\begin{itemize}
  \item Python code will follow PEP 8 guidelines, using pylint for linting.
  \item TypeScript code will follow the Microsoft Coding guideline, using Prettier for formatting and ESLint for linting.
  \item Git and GitHub will be used for version control, with feature branching and pull requests for code reviews.
  \item Unit tests, integration tests, and exploratory tests will be performed to ensure code works before merging.
\end{itemize}

\newpage{}

\section*{Appendix --- Reflection}

\begin{enumerate}
  \item \textbf{Why is it important to create a development plan prior to starting the
  project?}
  \\Having a development plan was crucial for our team because it gave us a shared understanding
  of our goals and responsibilities. Without a plan, we could easily duplicate work or miss deadlines,
  as we have seen in past projects where jumping straight into implementation led to confusion and rework.
  We wanted clear responsibilities, structured communication, and organized meeting roles from the start.
  We also felt that documenting agendas, workflows, and expectations would help us stay accountable and avoid
  wasting time in unstructured discussions. Overall, the plan helped us coordinate efficiently, understand
  dependencies, and move forward with more confidence.
  \item \textbf{In your opinion, what are the advantages and disadvantages of using
  CI/CD?}
  \\We think one of the biggest advantages of CI/CD is that it helps catch integration issues early and reduces
  the risk of major problems appearing late in development. It also helps keep the codebase stable when multiple
  people are working on interconnected features at the same time. At the same time, we recognize that setting up
  CI/CD can be time consuming at the beginning, and debugging failing automated checks can sometimes be frustrating.
  Even with those challenges, we believe the long-term benefits, such as better reliability, smoother collaboration,
  and faster detection of errors, make CI/CD worthwhile for our project.
  \item \textbf{What disagreements did your group have in this deliverable, if any,
  and how did you resolve them?}
  \\We had a few meaningful disagreements while preparing this deliverable, mainly around the level of detail
  needed in our workflow and how to structure GitHub issues. Some of us preferred a more detailed process to avoid
  confusion, while others felt that too much structure could slow us down and make meetings less efficient. We resolved
  this by discussing our previous experiences, listening to each other\'s reasoning, and agreeing on a balanced approach.
  In the end, we documented the main workflow in the development plan while leaving flexibility for finer technical details
  within GitHub issues. This helped us create a process that supports both structure and adaptability.
  \end{enumerate}

\section*{Appendix --- Team Charter}

\wss{borrows from
\href{https://engineering.up.edu/industry_partnerships/files/team-charter.pdf}
{University of Portland Team Charter}}

\subsection*{External Goals}

The long term goal for our team is to eventually make Proxi as our software business. In the short term goal is to create 
a product that can be showcased at the Capstone EXPO that can demonstrate not only technical functionality but also creativity, 
innovation, and accessibility. This project can be highlighted in future interviews, showcasing our collaboration and problem-solving 
skills and hopefully land us jobs. Another short term goal is that we are striving to receive an A+ in the course by delivering a high-quality, 
well documented, and fully integrated project.This dual focus on excellence and potential entrepreneurship motivates the team to deliver the 
project at the highest standard.

\subsection*{Attendance}

\subsubsection*{Expectations}

Expectations are clearly communicated to all team members regarding attendance and participation in meetings. So all team members are on the same page,
that we want at least 90 percent attendance rate for all scheduled meetings. Each team member is expected to be punctual and fully engaged during meetings that 
will run for 45 minutes each week as discussed. As a full time student we know it might not be possible to always meet in person, so virtual attendance is acceptable 
if a member cannot make it in person. We do agreed that in case we are meeting virtually, we will have our cameras on to ensure active participation. 
If a member cannot attend a meeting, they are expected to notify the team at least 24 hours in advance either through Discord or email. If a member misses 
more than one meeting without a valid excuse, the team will discuss the situation and may involve the TA if necessary.

\subsubsection*{Acceptable Excuse}
\begin{itemize}
  \item Illness or medical emergencies
  \item Family emergencies
  \item Unavoidable class/lab conflicts
  \item Scheduled exams
  \item Transportation issues (with advance notice)
  \item Technical problems (e.g., laptop malfunction)
\end{itemize}

\subsubsection*{Unacceptable Excuse}
\begin{itemize}
  \item Forgetting about meetings or deadlines
  \item Personal errands during scheduled team time
  \item Last-minute cancellations without valid reason
  \item Preventable conflicts that weren't communicated in advance
\end{itemize}

All absences require timely communication to the team. The team expects honesty and responsibility 
from all members. When valid excuses arise, the team will work together to redistribute tasks and 
maintain project momentum.

\subsubsection*{In Case of Emergency}

In case of an emergency and a team member cannot attend a scheduled meeting or meet a deadline, they should 
notify the team as soon as possible. Ideally notice should be at least 6 to 8 hours before the meeting or deadline, explaining 
the situation and providing an estimated time to catch up. The team will temporarily adjust responsibilities to 
keep the project on track. A member may miss no more than one meeting or one major deadline due to an emergency without 
extra discussion; missing more than that requires a team plan to catch up and may involve adjusting responsibilities or 
notifying the TA. Once the emergency is resolved, the member must promptly catch up on all missed work, review meeting notes, 
and resume contributions to upcoming tasks and deadlines, ensuring the team maintains momentum and uses scheduled time effectively.

\subsection*{Accountability and Teamwork}

\subsubsection*{Quality} 

It is expected that all members ome to meetings fully prepared may it be having reviewed their assigned tasks or notes, 
or any relevant materials ahead of time. Deliverables brought to the team whether code, documentation, test cases, 
or research should meet a high standard of clarity, completeness, and correctness. Before any work is submitted or 
integrated into the project, it will be reviewed by at least one other team member to catch errors, improve clarity, 
and ensure consistency. Work should be well-organized and easy for other team members to understand. By maintaining 
high-quality deliverables, thorough preparation, and a peer review process, the team can avoid technical debt, ensure 
smooth collaboration, and create a professional product that not only meets academic goals like earning an A+ but also 
has the potential to be refined and sold as a market-ready solution in the future.

\subsubsection*{Attitude}

Everyone from our team is expected to be respectful, cooperative, friendly. All the team members should be able to freely communicate their ideas. 
We want meetings to be friendly and positive and it is expected to actively participate in discussions, listen to different perspectives, and give 
helpful feedback when reviewing each other’s work and avoid negative critical tone. Our team do believe that a positive attitude is important because it not only 
boosts individual performance but also helps the whole team collaborate more effectively and stay motivated. 
We know that in a team conflicts are bound to come up and if they do, we will first talk it through as a team. This make sures that everyone has a chance 
to explain their side and present proofs of why they think one solution is better than another if its a conflict based on different solutions proposed.

\subsubsection*{Stay on Track}

To keep the team on track, we’ll rely on GitHub projects and issues to organize tasks, 
track deadlines, and see everyone’s progress in real time. We’ll review completed work, 
upcoming deadlines, and any blockers at every weekly 45-minute meeting so nothing falls through 
the cracks. Everyone is expected to attend most meetings (90 percent benchmark), 
complete their assigned tasks on time,and have their work peer-reviewed before submission. Team members who consistently do a great job 
will be recognized during meetings, and small rewards like leading a discussion or picking a fun 
team activity can be given. If someone falls behind or misses targets, the team will first check 
in with them to see if they need support and may redistribute tasks if needed. Repeated issues might 
involve minor consequences, like bringing snacks or coffee to the next meeting for the highest contributor on the project, or in serious cases, 
asking the TA for guidance. Due to this clear tracking, regular check-ins during our meets, and open communication, we make 
sure everyone contributes fairly. It will also help the team to stay coordinated, and the project keeps moving forward 
smoothly toward both our A+ goal and the long-term vision of creating a market-ready product.

\subsubsection*{Team Building}

We’ll be working together for a full year, so our team wants to make every meeting and work session something 
we actually look forward to. At the start of each meeting, we might do a quick check in where everyone shares 
a fun story, a small win from the week, or even a random interesting fact just enough to get everyone laughing 
or thinking. When we hit milestones, example finishing a big deliverable or meeting a tight deadline, we will
celebrate with team dinners or casual outings. This will help us to enjoy the moment and connect outside of coding. We also 
plan to have spontaneous brainstorming sessions, mini challenges, or even 10-minute games to keep our energy up and 
creativity flowing. Spending a whole year together gives us a rare chance to really learn each other’s strengths, 
keep motivation high, and grow Proxi beyond a school project, this is something we hope could eventually turn 
into a real business. By keeping things lively, flexible, and supportive, we make sure working on this project is fun, 
productive, and memorable for the whole team.

\subsubsection*{Decision Making} 

The majority of the decisions made by our team have be and will always be discussed collectively so that everyone can contribute their thoughts.  
We quickly vote to determine the best course of action if we are unable to agree either through Discord or 
by show of hands in person.  When differences arise, we make sure that everyone listens to one another and 
shares their point of view.  We seek the TA for assistance or take a brief break if we are still unable to make up our minds.  
Everyone will feel heard in this way, and we can continue to work efficiently and remain inspired on Proxi.

\end{document}